\documentclass[letterpaper]{article}
\usepackage[submission]{aaai2027}
\usepackage{times,helvet,courier}
\usepackage[hyphens]{url}
\usepackage{graphicx}
\usepackage{natbib}
\usepackage{caption}
\usepackage{booktabs,amsmath,amssymb,array,xspace}
\frenchspacing
\urlstyle{rm}
\def\UrlFont{\rm}
\pdfinfo{/TemplateVersion (2027.1)}
\setcounter{secnumdepth}{2}
\newcommand{\method}{TRACE\xspace}
\newcommand{\sg}{\operatorname{sg}}
\newcommand{\pending}{\textit{pending}}
\newcommand{\mainref}[1]{\ref{#1}}

\title{TRACE: Context-Probed Region Allocation for Test-Time Adaptation\\in Open-Vocabulary Semantic Segmentation}
\author{Anonymous Authors}
\affiliations{}
\begin{document}
\raggedbottom
\maketitle

\begin{abstract}
Open-vocabulary semantic segmentation assigns text-specified categories to image locations, while test-time adaptation learns from these predictions without labels. A confident local response does not reveal whether it remains supported when its surroundings change. We introduce TRACE, which uses a candidate-preserving context probe to allocate self-supervision. Connected predicted regions define the diagnostic units. For each region, one transformed input preserves its pixels and modifies its surroundings; the original query's response at the same locations supplies a fixed, detached support weight. TRACE normalizes these weights within each source image and applies them to region-averaged, prompt-specific entropy. This separates regional allocation from a direct pixel-area multiplier and gives each active image the same normalized loss budget. Only visual LayerNorm affine parameters adapt. Across six complete evaluation protocols, TRACE improves over matched initial-confidence weighting, with gains of 1.77 mIoU on COCO Object and 7.32 on COCO-CD. Explicit zero-, single-, and dual-probe controls examine the allocation and computational tradeoffs.
\end{abstract}

\section{Introduction}
Open-vocabulary semantic segmentation (OVSS) predicts a category at each image location from a supplied text vocabulary. Vision--language models provide transferable image--text representations, and dense prediction methods make them spatially accessible~\cite{radford2021clip,hajimiri2025naclip}. Test-time adaptation (TTA) adds a different decision: which unlabeled predictions should influence the model's parameters? Entropy-based adaptation sharpens current predictions~\cite{wang2021tent,noori2025mlmp}, making the choice of self-supervision important even when the initial segmentation appears confident.

Local confidence describes the model's response under the observed image, not its dependence on the surrounding content. This distinction matters in OVSS because local text matching is built from contextualized features learned largely through image--text alignment. Patch-aligned learning, local self-distillation, and modified token interactions address the global-to-local transfer~\cite{mukhoti2023pacl,wu2024clipself,shao2024cliptrase}. Object--context and region--language studies further examine how neighboring semantics affect dense representations~\cite{kim2025cass,li2025densevlm}. These methods improve the representations used for prediction; they do not by themselves decide which test-time region should contribute more to a parameter update.

Context need not be removed from prediction to make this distinction useful. Surrounding information can aid recognition~\cite{xiao2021backgrounds}; the additional question is how a candidate's response behaves when that information changes. Figure~\ref{fig:teaser} illustrates this question with actual model responses. Keeping a candidate's own pixels unchanged, we modify only its surroundings and read the same query at the same locations. Initially confident regions can then receive different support scores. This is a behavioral measurement under a specified intervention, rather than identification of a unique causal source of confidence.

\begin{figure}[t]
\centering
\includegraphics[width=\columnwidth]{figure/image_full_teaser.pdf}
\caption{Initial confidence and support after a candidate-preserving context change. Candidate pixels and diagnostic locations remain fixed. Labels show stored initial confidence and measured context-probe support from the evaluated single-probe method. These examples illustrate different responses, not their population frequency; Appendix~\ref{app:figures} describes selection and rendering.}
\label{fig:teaser}
\end{figure}

We use this measurement for \emph{regional allocation}, not for deleting contextual evidence from the final image. TRACE constructs connected regions from initial token predictions, bounds their number by top-$K$ confidence per source image, and evaluates one context probe per selected region. The detached score weights that region's average entropy. We normalize regional weights within each image before averaging image objectives. Consequently, neither a larger region nor an image with more total candidate support automatically receives a larger normalized loss budget.

Across six complete category protocols, the single-probe method improves on matched initial-confidence weighting in every protocol. The largest differences are on COCO-CD and COCO Object, at +7.32 and +1.77 mIoU; the COCO-Stuff difference is +0.04. Separate zero-, single-, and dual-probe comparisons on fixed training-source panels examine probe choice and cost, with a global mode selected on calibration before confirmation. Published MLMP values provide external operating-point references, separate from attribution within the same candidate and optimization pipeline.

Our contributions are:
\begin{itemize}
\item A candidate-preserving context probe that supplies fixed regional support for OVSS test-time self-supervision using one diagnostic input per candidate.
\item A source-image-normalized, region-averaged entropy objective that separates relative support from explicit area and cross-image normalization effects.
\item Six full-split matched allocation comparisons, calibration-locked probe comparisons under three fixed input orderings, and actual forward-count and timing measurements.
\end{itemize}

\section{Related Work}
\paragraph{Open-vocabulary dense prediction.}
OVSS transfers vision--language representations through dense matching, mask classification, or segmentation-oriented adaptation~\cite{ghiasi2022openseg,liang2023ovseg,cho2024catseg}. PACL aligns local patches with language, while CLIPSelf distills local visual features~\cite{mukhoti2023pacl,wu2024clipself}. NACLIP and CLIPtrase modify or analyze spatial-token interactions for training-free segmentation~\cite{hajimiri2025naclip,shao2024cliptrase}. TRACE addresses the subsequent allocation of test-time self-supervision, using the dense predictor as its inference interface.

\paragraph{Test-time adaptation.}
TENT optimizes entropy through normalization-affine parameters~\cite{wang2021tent}. CoTTA, EATA, and SAR introduce consistency, selection, or optimization controls~\cite{wang2022cotta,niu2022eata,niu2023sar}. TPT and WATT adapt vision--language models through prompts or template-associated weights~\cite{shu2022tpt,osowiechi2024watt}; Seg-TTO and SemLA address segmentation-specific adaptation~\cite{desilva2025segtto,qorbani2025semla}. MLMP combines visual levels and prompt-specific entropy with a global image-level term~\cite{noori2025mlmp}. TRACE retains dense prompt-specific entropy but assigns its influence through candidate-context responses and within-image regional normalization.

\paragraph{Context and model probing.}
Background-dependence studies distinguish recognition from the evidence sustaining it~\cite{xiao2021backgrounds}. CASS models object--context relations in OVSS, while DenseVLM examines biases in region--language alignment~\cite{kim2025cass,li2025densevlm}. Perturbation-based explanations inspect model behavior under controlled input changes~\cite{fong2019extremal,zhang2020conta}. TRACE uses a fixed, candidate-preserving perturbation to obtain relative update weights; it does not train an explanation model or assume calibrated probabilities of correctness.

\section{Method}
\label{sec:method}
An episode starts from reset model parameters $\theta_0$ and unlabeled source images $X_j$. TRACE selects connected predicted regions, measures their support under a candidate-preserving context change, and uses fixed regional weights to adapt the visual encoder. Figure~\ref{fig:pipeline} summarizes these three steps. Diagnostics never use image annotations.

\begin{figure*}[t]
\centering
\includegraphics[width=\textwidth]{figure/image_full_pipeline.pdf}
\caption{TRACE with one context probe per candidate. Initial token predictions define connected regions and source-image top-$K$ selection. Each candidate keeps its pixels while its surroundings undergo the fixed transform. The reset encoder measures the candidate query at its original token locations. Detached scores are normalized within source images and weight regional prompt entropy on unmodified inputs. Only visual LayerNorm affine parameters update; the inference readout reconstructs the final segmentation.}
\label{fig:pipeline}
\end{figure*}

\subsection{Candidate Region Construction}
Regions provide a spatial unit between an individual token and an entire image. They can be constructed from predictions without instance annotations or a separate proposal network. Here a region means a connected prediction component, not an estimated physical object instance.

Each source image is preprocessed into normalized input crops $x_b$, with source index $j(b)$. A $224\times224$ crop contains $16\times16$ spatial tokens for ViT-L/14; a ViT patch is distinct from an input crop. Let $\mathcal C$ be the text-query vocabulary, including any extensions. It need not equal the final evaluation-category set $\mathcal Y$. The reset diagnostic branch defines
\begin{equation}
\pi^0_{b,i}(c\mid x_b)=\operatorname{softmax}_{c\in\mathcal C}
z^{\mathrm{eval}}_{\theta_0}(x_b)_{i,c}.
\label{eq:reference}
\end{equation}
This branch uses normalized mean prompt embeddings and entropy-weighted visual-layer aggregation. Gradient-bearing optimization uses separate prompt distributions; Appendix~\ref{app:readout} specifies both readouts.

For each crop, we take the query argmax and extract four-neighbor connected components within its valid token mask. A candidate $R=(b_R,\Omega_R,c_R)$ records the crop, token locations, and predicted query. Its initial confidence is $\kappa_R=|\Omega_R|^{-1}\sum_{i\in\Omega_R}\pi^0_{b_R,i}(c_R\mid x_{b_R})$. Proposals from all crops of the same source image compete for its top-$K$ positions by descending $\kappa_R$; deterministic tie rules appear in Appendix~\ref{app:protocol}. The selected set for image $j$ is $\mathcal R_j$.

Confidence bounds the diagnostic workload and identifies predictions likely to be used for self-supervision; it does not certify reliability. The next module assigns weights independently of the initial rank. Components are not stitched across crops or merged across query extensions, even when extensions map to the same final category. Candidate geometry and query identities remain fixed within the episode.

\subsection{Candidate-Preserving Context Probe}
The probe asks whether the candidate's predicted query remains supported after a specified change to its surroundings. Preserving candidate pixels avoids deleting the semantic content whose response is measured. All categories use the same construction.

Let $u=\operatorname{Denorm}(x_{b_R})$ be the crop in RGB pixel space, and $M_R$ its token mask upsampled by nearest-neighbor interpolation to the crop dimensions. For the crop's per-channel spatial mean $\mu$, define
\begin{equation}
A(u)=\operatorname{clip}_{[0,255]}
\left(\operatorname{Shift}\bigl(\mu+0.35(u-\mu)\bigr)\right),
\end{equation}
where $\operatorname{Shift}(r,g,b)=(g,b,r)$. The diagnostic input is
\begin{equation}
x_R^{\mathrm{ctx}}=\operatorname{Norm}
\left(M_R\odot u+(1-M_R)\odot A(u)\right).
\label{eq:context-view}
\end{equation}
Only the mask complement is replaced. The response uses the original query and token locations:
\begin{equation}
\rho_R=\frac{1}{|\Omega_R|}\sum_{i\in\Omega_R}
\pi^0_{b_R,i}(c_R\mid x_R^{\mathrm{ctx}}).
\label{eq:support}
\end{equation}
This is a mean of probabilities, not logits. No new region is extracted from the transformed prediction. The score measures support under this context change, not invariance to every possible context. We call this probe mode $C$ in the comparisons below.

Each region receives its own transformed input, but regions do not require serial model calls. We index views by candidate, stack them into bounded batches, and map their logits back to the corresponding candidate before averaging over $\Omega_R$. Diagnostic inputs therefore number one per candidate, while encoder calls depend on view batch size. The resulting scores are detached and fixed during the episode.

\subsection{Image-Balanced Regional Adaptation}
Support and size should play distinct roles in allocating self-supervision. Normalizing scores across an entire episode also makes images compete through a shared denominator. TRACE instead normalizes within each source image, then averages the active image objectives.

Let $p^a_{\theta,b,i}$ be the current query distribution for prompt $a\in\{1,\ldots,J\}$ on an unmodified crop, and $H(p)=-\sum_{c\in\mathcal C}p(c)\log p(c)$. Each candidate has mean prompt/token entropy
\begin{equation}
\overline H_R(\theta)=\frac{1}{J|\Omega_R|}
\sum_{a=1}^{J}\sum_{i\in\Omega_R}H\bigl(p^a_{\theta,b_R,i}\bigr).
\end{equation}
Define active regions $\mathcal R_j^+=\{R\in\mathcal R_j:\rho_R>\epsilon\}$, with the fixed numerical threshold $\epsilon=10^{-8}$, and active images $\mathcal I_+=\{j:\mathcal R_j^+\neq\varnothing\}$. We optimize
\begin{align}
\alpha_{R\mid j}&=\frac{\sg(\rho_R)}
{\sum_{Q\in\mathcal R_j^+}\sg(\rho_Q)},\\
\mathcal L_{\mathrm{TRACE}}(\theta)&=\frac{1}{|\mathcal I_+|}\sum_{j\in\mathcal I_+}
\sum_{R\in\mathcal R_j^+}\alpha_{R\mid j}\overline H_R(\theta).
\label{eq:loss}
\end{align}
Here $\sg$ denotes stop-gradient. Images without active candidates contribute no term; if no image is active, adaptation is skipped. The candidate query selects the diagnostic response but is not a cross-entropy pseudo-label for this entropy objective.

Each active image has total coefficient mass $1/|\mathcal I_+|$. Within it, $\alpha_{R\mid j}$ gives relative regional priority without an explicit area multiplier. This is a loss-budget property, not equality of gradient norms or invariance to batch composition. Common positive scaling of an image's active scores cancels; a singleton candidate's positive score also cancels. Support thus controls relative allocation, not the absolute learning rate or number of updates (Appendix~\ref{app:theory}).

Only visual LayerNorm affine parameters update using Adam; the text encoder and other visual parameters remain frozen. Model and optimizer reset before each episode. Final inference uses unmodified inputs, reconstructs pixel query scores, and takes the maximum over extensions of each evaluation category before category argmax. Appendix~\ref{app:protocol} gives this coordinate and vocabulary mapping.

\section{Experiments}
\subsection{Evaluation Protocol}
We ask whether probing improves allocation beyond initial confidence, whether one view is a useful operating point, and what computation it requires. Evaluation covers six complete protocols: VOC20/VOC21 (1,449 images each), PC60 (5,105), COCO Object (5,000), COCO-CD (500), and COCO-Stuff (5,000). VOC20/VOC21 share source images but differ in categories. We pool per-class intersection and union over each complete split, then average IoU over positive-union classes.

All local arms use NACLIP ViT-L/14, the last 18 visual outputs, seven prompts, and the same checkpoint. Source images and evaluation labels are resized to $224\times224$, using nearest-neighbor interpolation for labels. The crop size is also $224\times224$, so these measurements use one crop per resized image, not native-resolution evaluation. Uniform, Confidence, and TRACE use identical candidates, order, readouts, learning rate, update steps, and episodic initialization. Each arm executes independently and resets its optimizer. All six splits' predictions are sealed before offline scoring.

The global probe mode is selected using two fixed 128-image training calibration panels, then frozen before 256-image confirmation panels under three predeclared orderings. This stage performs no full-evaluation mode or parameter search. The inherited configurations and official evaluation splits have prior development exposure. Training-panel exclusions cover documented historical selections, but incomplete July candidate-pool mappings prevent certification that all historical exposure was excluded. We distinguish calibration, descriptive confirmation, and development-exposed full evaluation, rather than presenting them as globally unseen validation. Appendix~\ref{app:selection} gives the selection rule and complete outcomes.

\subsection{Full-Split Regional Allocation}
\label{sec:allocation}
Matched controls isolate the information supplied by the probe. Uniform assigns equal regional weights; Confidence uses initial $\kappa_R$; TRACE uses the context-probe score. All three normalize within images and average regional entropy. Only TRACE constructs diagnostic views.

\begin{table}[t]
\caption{Seed-0 full official splits with development exposure and inherited per-dataset configurations; fixed calibration-selected C, no reselection. Descriptive results, not unseen/equal-tuning or significance evidence. Pooled class mIoU includes every image, including episode 0. Only newly executed matched controls. Object/Stuff denote COCO Object/COCO-Stuff. $\Delta$ is TRACE(C) minus Confidence in percentage points.}
\label{tab:image-full-comparison}
\centering\small\setlength{\tabcolsep}{4pt}
\begin{tabular}{@{}lrrrr@{}}
\toprule
Dataset & Uniform & Confidence & TRACE(C) & $\Delta$ \\
\midrule
PC60 & 30.38 & 30.34 & 30.69 & +0.36 \\
Object & 29.81 & 30.83 & 32.59 & +1.77 \\
VOC21 & 52.44 & 52.33 & 52.62 & +0.30 \\
VOC20 & 81.78 & 82.18 & 83.49 & +1.31 \\
COCO-CD & 66.22 & 67.66 & 74.98 & +7.32 \\
Stuff & 21.74 & 21.79 & 21.84 & +0.04 \\
\bottomrule
\end{tabular}
\end{table}


TRACE improves over Confidence on all six protocols (Table~\ref{tab:image-full-comparison}): +1.31 on VOC20, +0.30 on VOC21, +0.36 on PC60, +1.77 on Object, +7.32 on CD, and +0.04 on Stuff. Its largest differences from Uniform are +8.77 on CD and +2.79 on Object. Fixed candidates and optimizer settings make these differences evidence for allocation rather than a different proposal network or update budget.

The effect sizes identify where this information is useful. On CD, initial confidence improves Uniform by 1.45 points, while the additional context response contributes a further 7.32. On Object, the corresponding differences are 1.02 and 1.77. In these conditions, selecting which candidate receives weight matters beyond preferring confident proposals. On Stuff, Confidence and TRACE are within 0.05 points: the evidence supports compatibility with this mixed-category protocol, not a large independent effect there. These controls do not isolate candidate construction or calibrate individual update utility.

\begin{figure}[t]
\centering
\includegraphics[width=\columnwidth]{figure/image_full_allocation.pdf}
\caption{TRACE minus matched Confidence and Uniform allocation, in mIoU percentage points. All six protocols use the same candidate and optimization pipeline within each comparison. Policies differ in allocation scores; TRACE alone adds a context probe.}
\label{fig:allocation}
\end{figure}

\subsection{Published and Local References}
Table~\ref{tab:image-full-reference} reproduces the clean NACLIP values reported by MLMP~\cite{noori2025mlmp}, including reported dispersion, and separately lists our local single-seed results. TRACE reaches 83.49/52.62/30.69/32.59/21.84 on VOC20/VOC21/PC60/Object/Stuff. Against the strongest reported method in each column, MLMP, the differences are $-0.27/ +1.84/ +2.70/ +3.75/ +0.59$ points.

\begin{table*}[t]
\caption{R: means and reported standard deviations from main Table~5 (Original NACLIP rows) of \citet{noori2025mlmp}, freshly cross-checked against the source PDF; not rerun. No Adapt has no reported standard deviation. L: newly measured TRACE(C), one deterministic seed, no standard deviation. R and L are reference operating points, not matched checkpoint/loader/query/episode/budget/tuning protocols. The negative VOC20 difference is retained.}
\label{tab:image-full-reference}
\centering\small
\begin{tabular}{@{}lrrrrr@{}}
\toprule
Method & VOC20 & VOC21 & PC60 & Object & Stuff \\
\midrule
No Adapt (R) & $75.91$ & $45.12$ & $24.95$ & $23.80$ & $18.34$ \\
TENT (R) & $77.00 \pm 0.04$ & $45.65 \pm 0.02$ & $25.29 \pm 0.01$ & $24.88 \pm 0.01$ & $18.76 \pm 0.01$ \\
TPT (R) & $75.93 \pm 0.01$ & $45.17 \pm 0.01$ & $24.98 \pm 0.01$ & $23.84 \pm 0.01$ & $18.35 \pm 0.01$ \\
WATT (R) & $57.73 \pm 0.06$ & $28.58 \pm 0.05$ & $14.77 \pm 0.03$ & $14.14 \pm 0.06$ & $9.49 \pm 0.02$ \\
CLIPArTT (R) & $72.77 \pm 0.14$ & $39.50 \pm 0.04$ & $21.88 \pm 0.03$ & $21.34 \pm 0.03$ & $15.48 \pm 0.01$ \\
MLMP (R) & $83.76 \pm 0.00$ & $50.78 \pm 0.02$ & $27.99 \pm 0.03$ & $28.84 \pm 0.01$ & $21.25 \pm 0.01$ \\
\midrule
TRACE(C) (L) & 83.49 & 52.62 & 30.69 & 32.59 & 21.84 \\
$\Delta$ vs. best R (pp) & -0.27 & +1.84 & +2.70 & +3.75 & +0.59 \\
\bottomrule
\end{tabular}
\end{table*}
\begin{table*}[t]
\caption{COCO-CD local reference operating points. Historical native baselines are sealed existing measurements, not rerun or part of the new matched three-arm experiment. Checkpoint and full input order match recorded metadata; method backends, prompts and optimization budgets retain their own protocols. TENT is the best historical native baseline. The historical No Adapt row is not a certified reset readout of the new method.}
\label{tab:image-full-cd-reference}
\centering\small
\begin{tabular}{@{}lr@{}}
\toprule
Method / protocol & COCO-CD \\
\midrule
No Adapt (L, historical native) & 61.14 \\
TENT (L, historical native) & 72.20 \\
TPT (L, historical native) & 64.27 \\
MLMP (L, historical native) & 71.01 \\
\midrule
TRACE(C) (L, new) & 74.98 \\
$\Delta$ vs. native TENT (pp) & +2.78 \\
\bottomrule
\end{tabular}
\end{table*}


Published methods retain their own readouts and settings. A shared backbone name does not establish identical checkpoints, extensions, evaluation geometry, or optimization budgets. In particular, the VOC20 reference exceeds our score. We do not interpret cross-paper differences as matched superiority; the allocation controls provide direct attribution within our implementation.

COCO-CD has no entry in that published table. Its completed local native comparison gives 72.20 mIoU for the strongest baseline, TENT, versus 74.98 for selected TRACE, a difference of 2.78 points. Native methods share the ordered CD panel but retain their method-specific backends and settings (Appendix~\ref{app:native}). This is a separate operating-point comparison, not an equal-budget causal estimate.

\subsection{Probe Choice and Input Ordering}
Probe comparisons separate the value of probing from the cost of collecting multiple responses. We compare a candidate-preserving grayscale-blur view ($T$), the context transform ($C$), and their minimum under the image-normalized objective. Uniform and Confidence are zero-probe controls; an episode-normalized minimum separately tests the cross-image denominator. Candidate sets, configurations, and all six arms are fixed before prediction.

Calibration selects one global mode by its worst-dataset margin against image-normalized Confidence on PC60 and Object, with a fixed cost-aware tie rule. It selects $C$, sealed before confirmation prediction. The calibration PC60 margin is $-0.27$, so selection does not mean both calibration datasets improve. Appendix~\ref{app:selection} retains all outcomes rather than substituting per-dataset winners.

On each 256-image confirmation panel, selected $C$ exceeds Confidence under all three orderings: +0.07--0.18 on PC60 and +2.10--2.92 on Object. $T$ is higher on Object; the original episode-normalized minimum is higher on PC60. These results support a single-probe operating point, not universal optimality of $C$ or image normalization. The latter removes an explicit cross-image score denominator algebraically, but shared gradients and Adam still make grouping relevant. Reordering the same images diagnoses composition dependence; it is not independent population sampling or a confidence interval.

\subsection{Actual Computation}
One context probe requires $N_{\mathrm{cand}}$ diagnostic inputs, versus $2N_{\mathrm{cand}}$ for a dual-probe minimum. With $P$ input crops and $S$ update steps, a nonempty episode uses
\begin{equation}
N_{\mathrm{enc\ inputs}}=N_{\mathrm{cand}}+(S+2)P,
\end{equation}
including initial readout, adaptation, and final inference. Zero-probe controls omit the first term. Diagnostic \emph{calls} depend on the view batch size; input count is not a wall-time speedup.

Full-split independent-arm measurements give 119.41 ms/image for TRACE versus 83.73 for Confidence on PC60, and 292.00 versus 237.41 on Object, alongside +0.36/+1.77 mIoU. On Stuff, 186.14 versus 129.35 ms/image accompanies a +0.04-point difference. The measured return therefore depends on how much allocation changes the task outcome. Appendix~\ref{app:cost} reports every arm, actual input/call counts, and timing denominators.

These are synchronized model-side measurements including reset, diagnosis, optimization, reconstruction, CPU prediction transfer, and finite-value observations. They exclude setup, decoding, shared transfer/geometry preparation, artifact writes, and scoring. Only episode 0 is omitted from time; all images remain in accuracy and work counts. Earlier training-panel single/dual-probe timings use their own instrumentation; their percentage differences are not transferred to this full stage.

\subsection{Qualitative Allocation Effects}
Figure~\ref{fig:segmentation-cases} compares actual Confidence and TRACE segmentations from the new full evaluation. The comparator changes allocation, not the candidate generator or initialization, making the examples inspectable illustrations of the joint update. Positive and negative examples are retained. The selected class-level IoUs concern the displayed images, not full-dataset performance.

\begin{figure*}[t]
\centering
\includegraphics[width=\textwidth]{figure/image_full_segmentation.pdf}
\caption{Actual selected segmentations: Input / Confidence / TRACE / Ground truth. Confidence is the matched initial-confidence allocation control, not MLMP. Annotations are IoU for the named class in the displayed image, not dataset mIoU. Cases include positive and negative differences and are selected for diagnostic visibility, not representativeness. GT is used only offline. Appendix~\ref{app:figures} gives identities and selection rules.}
\label{fig:segmentation-cases}
\end{figure*}

\section{Discussion and Scope}
TRACE measures support under a particular context change and uses it for relative allocation. High support can coexist with shared errors or sensitivity to untested transformations. Connected regions can contain mixed-quality positions, and top-$K$ can exclude useful low-confidence candidates. Region averaging and equal image budgets constrain coefficients, not gradient magnitudes or an individual update's realized value.

The evidence directly supports allocation effects in the six matched protocols and the specified computation. Its scope is one backbone, fixed transformation, inherited per-dataset settings, and resized evaluation geometry. Development exposure and published-reference protocol differences limit claims about untouched generalization and matched external competitiveness. These boundaries do not alter the fixed-candidate comparison, but prevent interpreting it as universal update safety or state-of-the-art performance under every protocol.

\section{Conclusion}
TRACE turns one candidate-preserving context response into a fixed regional allocation score for OVSS test-time adaptation. Normalizing within source images and averaging regional entropy makes the loss budget explicit while retaining the dense prediction pipeline. Full-split matched comparisons show that the additional response can improve on initial-confidence weighting, especially on COCO-CD and COCO Object. Calibration-locked probe comparisons and actual execution measurements establish the corresponding single-probe operating point and its computational tradeoff.

\section*{AI Assistance}
AI assistance was used for source drafting, experimental orchestration, mathematical consistency checks, visualizations, and table-provenance tooling. The authors remain responsible for verifying the experimental evidence and final text.

\bibliography{references}
\clearpage
\appendix
\section{Implementation and Reproduction Protocol}
\label{app:protocol}
\subsection{Method Identity and Execution}
The evaluated method is \texttt{ImageBalancedTRACE} with the single global mode \texttt{C}. It measures one candidate-preserving context-perturbed response, normalizes detached region weights within each source image, and averages the resulting objectives over active source images. The implementation is \path{trace_ovss/method/image_balanced.py}; the full-split worker is \path{tools/evolution_image_full/run.py}, using the inherited model, candidate, optimization, and reconstruction code. This is a different objective from episode-normalized two-view TRACE. It has no second probe in the selected mode, semantic-role branch, candidate deletion, dominance attenuation, fragment guard, or CLS loss.

An episode starts from the pretrained parameters and a reset optimizer. The reset model constructs candidates and measures their supports once without gradients. Adaptation then reuses the fixed candidate locations and weights on unmodified crops, followed by final inference on those crops. Only visual LayerNorm affine parameters change. All three full-split arms execute independently from the same initial model and RNG state; their execution order alternates by episode. Uniform and Confidence use the same source-image-normalized objective, with weights one and initial confidence, respectively, and compute no diagnostic views.

The runtime uses deterministic PyTorch algorithms with errors on nondeterministic operations, \texttt{CUBLAS\_WORKSPACE\_CONFIG=:4096:8}, deterministic cuDNN, and disabled cuDNN benchmarking. Default attention dispatch remains enabled. Model seed 0 and the original complete evaluation order are fixed. Deterministic execution is not an independent replication or an uncertainty estimate.

\subsection{Images, Crops, Tokens, and Candidates}
A source image is the unit of objective normalization; a crop is an encoder input; a token is a spatial encoder output. The general loader extracts complete sliding-window crops of size $224\times224$ with stride 112, using row-major grid positions and no cross-crop component stitching. Crops retain their source-image IDs. The evaluated configuration first bilinearly resizes every RGB image to $224\times224$, so it yields one crop per source image. This one-crop setting is an evaluation configuration, not the definition of the multi-crop algorithm. ViT-L/14 uses $14\times14$-pixel patches, giving a $16\times16$ spatial-token grid for each crop; 14 is not the crop size.

Let $\mathcal C$ be the supplied text-query vocabulary, including enabled query extensions, and let $\mathcal Y$ be the evaluation-category set. These sets need not coincide. On crop $b$, the reset-model distribution is
\begin{equation}
\pi^0_{b,i}(c\mid x_b)=
\operatorname{softmax}_{c\in\mathcal C}
z^{\mathrm{eval}}_{\theta_0}(x_b)_{i,c}.
\end{equation}
Four-connected components of its token-wise query argmax define proposals $R=(b_R,\Omega_R,c_R)$, with crop index $b_R$, nonempty token set $\Omega_R$, and predicted query $c_R$. Let $j(b)$ map a crop to its source image, and write $j_R=j(b_R)$ for bookkeeping; this is not an additional candidate component. The initial confidence is
\begin{equation}
\kappa_R=\frac{1}{|\Omega_R|}\sum_{i\in\Omega_R}
\pi^0_{b_R,i}(c_R\mid x_{b_R}).
\end{equation}
Proposals from all crops of source image $j$ compete for its first $K$ positions, forming the selected set $\mathcal R_j$. Ranking uses descending $\kappa_R$, descending token area, ascending query index, and ascending component index; stable crop/component enumeration resolves remaining exact ties. There is no overlap union, cross-crop stitching, or minimum-area filter beyond nonemptiness. A valid-token mask excludes geometric padding, not ground-truth background. Distinct query extensions remain distinct components even if they map to the same category.

\subsection{Exact Context Transform and Support}
The C probe changes only the complement of a candidate's input-space mask. Denormalize crop $x_b$ to RGB pixels $u_b\in[0,255]^{3\times224\times224}$ using mean $(122.7709,116.7460,104.0937)$ and standard deviation $(68.5005,66.6322,70.3232)$. For channel $k$, let $\mu_{b,k}$ be its mean over the \emph{entire crop}, including candidate pixels. Define
\begin{align}
A(u_b)&=\operatorname{clip}_{[0,255]}
\left(\operatorname{Shift}\bigl(\mu_b+0.35(u_b-\mu_b)\bigr)\right),\\
\operatorname{Shift}(r,g,b)&=(g,b,r),\\
u_R^{\mathrm{ctx}}&=M_R\odot u_{b_R}
 +(1-M_R)\odot A(u_{b_R}),
\end{align}
where $M_R$ is nearest-neighbor upsampling of the original binary token mask. The view is renormalized before encoding. There is no random transform sampling or learned transform parameter.

The detached C weight is the mean probability of the original query at its original token locations:
\begin{equation}
\rho_R=\frac{1}{|\Omega_R|}
\sum_{i\in\Omega_R}\pi^0_{b_R,i}
\bigl(c_R\mid\operatorname{Norm}(u_R^{\mathrm{ctx}})\bigr).
\label{eq:app-C}
\end{equation}
The diagnostic branch returns token-grid logits without pixel interpolation. It does not re-extract components, match transformed masks, use a probability ratio, or measure a CLS response. Each selected candidate schedules exactly one diagnostic input; view batching changes calls, not that input count. The selected C branch never constructs the unused T view.

Candidate preservation is geometric, not causal isolation. Candidate pixels remain unchanged, but the full-crop channel means depend on them, and the encoder's contextual interactions can mix information inside and outside the mask. The code's compatibility name \texttt{counterfactual\_support} denotes this specified perturbed response, not an identified causal counterfactual or a calibrated probability of update correctness.

\subsection{Dense Readouts and Prompt Entropy}
\label{app:readout}
Diagnosis and final inference use entropy-weighted visual-layer aggregation and normalized mean prompt embeddings. Optimization instead uses uniform mean-layer aggregation and separate prompt embeddings. Both readouts use the last 18 visual outputs of reduced NACLIP ViT-L/14, blocks 7--24 indexed from one, with neighborhood standard deviation 5~\cite{hajimiri2025naclip}. The selected tensors are attention-branch outputs; the ordinary residual and MLP path propagates between blocks.

Let $h_{\ell,b,i}$ be the selected pre-projection output, $\psi_\theta(h)=\mathrm{LN}_{\mathrm{post},\theta}(h)W_{\mathrm{proj}}$, and a hat denote unit-length normalization. For prompt $a\in\{1,\ldots,J\}$ and query $c$, $t_{a,c}$ is the normalized frozen text embedding, and $\bar t_c=\widehat{J^{-1}\sum_at_{a,c}}$, with $J=7$. The multiplier $\lambda=\exp(\ell_{\mathrm{CLIP}})$ is the frozen pretrained logit scale, not a tuned diagnostic temperature. Using $H(p)=-\sum_cp(c)\log p(c)$, inference first computes per-layer query logits and crop-specific entropy weights:
\begin{align}
d_{\ell,b,i}(c)&=\lambda\left\langle
\widehat{\psi_\theta(h_{\ell,b,i})},\bar t_c\right\rangle,\\
v_{\ell,b}&=\frac1N\sum_iH\!\left(
\operatorname{softmax}_c d_{\ell,b,i}\right),\\
w_{\ell,b}&=\frac{\exp(-v_{\ell,b})}{\sum_m\exp(-v_{m,b})}.
\end{align}
Here $N=256$ spatial tokens and the CLS token is excluded. Raw layer outputs are then combined before post-normalization and projection:
\begin{equation}
z^{\mathrm{eval}}_{b,i}(c)=\lambda\left\langle
\widehat{\psi_\theta\!\left(\sum_\ell w_{\ell,b}h_{\ell,b,i}\right)},
\bar t_c\right\rangle.
\end{equation}
This is not an average of per-layer probabilities. Diagnosis uses $\theta_0$; final inference uses the adapted parameters. For optimization, $p^a_{\theta,b,i}=\operatorname{softmax}_c z^{a,\mathrm{adapt}}_{b,i}$, with gradient-bearing prompt-specific logits
\begin{equation}
z^{a,\mathrm{adapt}}_{b,i}(c)=\lambda\left\langle
\widehat{\psi_\theta\!\left(\frac1{18}\sum_\ell h_{\ell,b,i}\right)},
t_{a,c}\right\rangle.
\end{equation}
Thus diagnostic probabilities and adaptation probabilities are related readouts, not interchangeable tensors.

The seven fixed templates are:
\begin{quote}\small
\texttt{itap of a \{\}}\\
\texttt{a bad photo of the \{\}.}\\
\texttt{a origami \{\}.}\\
\texttt{a photo of the large \{\}.}\\
\texttt{a \{\} in a video game.}\\
\texttt{art of the \{\}.}\\
\texttt{a photo of the small \{\}.}
\end{quote}
Prompts and text parameters are frozen. The official OpenAI CLIP \texttt{ViT-L-14.pt} checkpoint~\cite{radford2021clip} has SHA-256 given by concatenating
\begin{quote}\small\ttfamily
b8cca3fd41ae0c99ba7e8951adf17d26\\
7cdb84cd88be6f7c2e0eca1737a03836
\end{quote}
Resource hashes and the exact executed parameter names are retained with the frozen protocols.

\subsection{Source-Image-Normalized Optimization}
Let $\mathcal R_j^+=\{R\in\mathcal R_j:|\Omega_R|>0,\ \rho_R>\epsilon\}$, with $\epsilon=10^{-8}$, and $\mathcal I_+=\{j:\mathcal R_j^+\ne\varnothing\}$. The region entropy and objective are
\begin{align}
\overline H_R(\theta)&=\frac{1}{J|\Omega_R|}
 \sum_{a=1}^J\sum_{i\in\Omega_R}H(p^a_{\theta,b_R,i}),\\
Z_j&=\sum_{Q\in\mathcal R_j^+}\sg(\rho_Q),\\
\alpha_{R\mid j}&=\frac{\sg(\rho_R)}{Z_j},\qquad R\in\mathcal R_j^+,\\
\mathcal L_{\mathrm{TRACE}}(\theta)
 &=\frac1{|\mathcal I_+|}\sum_{j\in\mathcal I_+}
 \sum_{R\in\mathcal R_j^+}\alpha_{R\mid j}\,\overline H_R(\theta).
\label{eq:app-image-loss}
\end{align}
The stop-gradient operator $\sg$ emphasizes that neither supports nor their normalizers are differentiated. Candidate indices, source IDs, and weights remain fixed throughout the episode. Entropy is averaged over prompts and tokens \emph{within each region}, not pooled over all candidate pixels. The original query defines the support measurement but is not a cross-entropy pseudo-label.

Only images with an active nonempty region enter the outer denominator. With none, the tensor helper returns the graph-connected zero $0\sum H$ and the engine skips adaptation. The implementation groups terms by source ID, converts detached scalar weights to the entropy dtype, and uses a double-precision comparison of each summed denominator against $\epsilon$, falling back to one when the denominator does not exceed it. Equation~\ref{eq:app-image-loss} describes the finite active-score case. Crops of one image must share their source ID; treating each crop as a separate image would implement a different objective.

Adam uses $\beta_1=0.9$, $\beta_2=0.999$, and zero weight decay. Trainable parameters include visual input/output LayerNorms and the normalization layers along the executed visual blocks. The final selected block's unused MLP normalization is frozen in the reduced architecture; earlier executed blocks remain differentiable even when their outputs are not selected for readout. Text encoding, projection weights, and the remaining backbone parameters stay frozen. No global image-level or CLS-token objective is added.

\begin{table*}[t]
\caption{Frozen per-dataset settings; all allocation and spatial-normalization controls retain their dataset configuration. Batch denotes source images per episode, $K$ candidate limit, $B_v$ diagnostic view-batch size, and Ext. query extensions. Complete preprocessing and source identities are retained in the generated provenance.}
\label{tab:configurations}
\centering\small
\begin{tabular}{@{}lrrrrrr@{}}
\toprule
Dataset & Batch & LR & Steps & $K$ & $B_v$ & Ext. \\
\midrule
COCO-CD & 28 & 0.0086 & 5 & 4 & 16 & No \\
PC60 & 9 & 0.00225 & 11 & 22 & 16 & Yes \\
VOC21 & 16 & 0.00063 & 36 & 33 & 16 & Yes \\
COCO Object & 14 & 0.003 & 29 & 34 & 16 & Yes \\
COCO-Stuff & 5 & 0.0034 & 12 & 33 & 16 & Yes \\
VOC20 & 1 & 0.001105 & 9 & 8 & 16 & Yes \\
\bottomrule
\end{tabular}
\end{table*}


Table~\ref{tab:configurations} specifies the exact inherited settings, not newly optimized C settings. All matched arms retain the split's batch size, learning rate, steps, candidate limit, diagnostic batch size, and query-extension flag. The full worker inherits the sealed resolved configurations rather than substituting a current configuration lookup. These settings are not equal budgets across datasets or external methods.

\subsection{Pixel Reconstruction and Category Evaluation}
Final reconstruction operates on query logits before category collapsing. For each crop, bilinearly upsample its adapted $16\times16$ query-logit grid to $224\times224$ with \texttt{align\_corners=False}. Place the crop logits at their sliding-window offsets, sum overlaps, and divide by the pixelwise coverage count (clamped below by one). The general implementation uses complete windows, does not add a forced terminal crop, and leaves uncovered positions at zero logits; the present one-crop square configuration covers every evaluated pixel. Crops from one source are reconstructed together, not scored as independent images.

After reconstruction, apply the pixelwise query softmax to obtain $s_c(u)$. If $e(c)$ maps an extension query to its evaluation category, prediction is
\begin{equation}
g_y(u)=\max_{c:e(c)=y}s_c(u),\qquad
\widehat Y(u)=\arg\max_{y\in\mathcal Y}g_y(u).
\end{equation}
This is a maximum after pixel softmax, not a sum of extension probabilities, a category softmax before reconstruction, or a category collapse before candidate construction. Without extensions, $e$ is the identity. Vocabulary files under \path{trace_ovss/resources/class_names/} define categories by nonempty lines and comma-separated extensions within a line: \texttt{cls\_voc20.txt}, \texttt{cls\_voc21.txt}, \texttt{cls\_context60.txt}, \texttt{cls\_coco\_object.txt}, \texttt{coco\_cd.txt}, and \texttt{cls\_coco\_stuff.txt}.

The complete evaluation populations are 1,449 images each for VOC20/VOC21, 5,105 for PC60, 5,000 each for COCO Object/Stuff, and 500 for COCO-CD. VOC20 and VOC21 use different category protocols on the same images. Labels are nearest-neighbor resized to the $224\times224$ evaluation grid, not evaluated at native annotation resolution. The category counts with positive union are 20, 21, 60, 81, 171, and 3 for VOC20, VOC21, PC60, Object, Stuff, and CD. VOC21, PC60, Object, and CD include explicit background; VOC20 excludes background and Stuff uses its 171 object/stuff categories. Ignore label 255 is excluded. Pooling intersections $I_y$ and unions $U_y$ over the complete split gives
\begin{equation}
\operatorname{mIoU}=\frac{100}{|\mathcal Y_+|}
\sum_{y\in\mathcal Y_+}\frac{I_y}{U_y},
\qquad \mathcal Y_+=\{y:U_y>0\}.
\end{equation}
This is not an average of per-image or candidate IoUs.

\subsection{Frozen Evidence and Label Separation}
The full evaluation is \path{trace_image_full_20261029_0333} under \path{/home/wwww/hzl/results/}. Its \texttt{SUITE\_PREPARED.json} binds ordered manifests, inherited configurations, the selected-mode record, checkpoint/resources, and frozen runtime sources. Six four-image RGB-only gates compare worker C against direct \texttt{ImageBalancedTRACE(view\_mode="C")}, with exact candidate identity, support/coefficient/loss/gradient tolerance $10^{-5}$, zero changed final-label pixels, and actual zero/zero/one-probe schedules. These small gates check implementation identity, not full-split accuracy.

All six complete three-arm prediction sets are sealed before any evaluation-label contents are read. The prediction loader opens RGB images only; separate CPU scoring pools the saved final predictions. The root's \texttt{ALL\_FULL\_PREDICTIONS\_SEALED.json}, per-cell prediction/scoring seals, six canonical \texttt{result.json} files, and \texttt{SUITE\_SEALED.json} bind the 18 arm results. The prepared identity is the concatenation
\begin{quote}\small\ttfamily
e0f7374042eadb65fbc8146de6dfcf35\\
a71151dc1c95c115b0d4cf9c8444edaa
\end{quote}
The executed source snapshot, rather than an unspecified live checkout, defines the method for every new full result.

These official splits and inherited configurations were exposed during earlier development. Historical training-subset searches and some full-validation-conditioned choices preceded this evaluation: notably VOC20 learning-rate retention and COCO-CD/Stuff learning-rate refinements used full-validation results. The source of VOC21's 36-step choice and the earlier VOC20 candidate-limit choice is not fully established. The parameter-origin record is \path{trace_parameter_provenance_20261029_0030/audit.json} under the same results directory. Freezing the new C evaluation does not erase this exposure; the full results are fixed-configuration, development-exposed evidence rather than untouched confirmation or evidence of equal tuning opportunity.

\section{Objective Algebra and Interpretation}
\label{app:theory}
\subsection{Equal Source Budgets}
Source-image normalization fixes each active image's total objective-coefficient mass. For a finite fixed active set, the within-image weights satisfy
\begin{equation}
\sum_{R\in\mathcal R_j^+}\alpha_{R\mid j}=1,\qquad
\sum_{R\in\mathcal R_j^+}\frac{\alpha_{R\mid j}}{|\mathcal I_+|}
=\frac1{|\mathcal I_+|}.
\end{equation}
In contrast, episode normalization assigns source image $j$ the mass $Z_j/\sum_{k\in\mathcal I_+}Z_k$. Equal image budgets therefore remove this explicit competition through a shared sum-of-support denominator. They do not equalize candidate counts, region areas, gradient norms, Adam updates, or segmentation accuracy. Relative weights and the number of selected regions still affect allocation within an image.

\subsection{Rescaling and Singleton Limits}
Within-image positive rescaling cancels only while the active set and numerical behavior remain unchanged. Replacing every $\rho_R$ in source image $j$ by $\gamma_j\rho_R$, $\gamma_j>0$, leaves $\alpha_{R\mid j}$ and the objective unchanged under those conditions. If an image has exactly one active region, its within-image weight is one and its full objective coefficient is $1/|\mathcal I_+|$, regardless of that region's positive support; other images can still have multiple active regions. With a single active source image, the objective equals the episode-normalized objective for the same candidates and weights. These identities concern normalization, not empirical equivalence between a single probe and two probes.

A low absolute support does not by itself suppress an active image's normalized contribution. Threshold crossings change the active set, and changing the number of active images changes the outer factor. Thus the C score supplies relative allocation within an image, not an absolute rejection or update-safety guarantee.

\subsection{Gradient and Area Boundaries}
Detached coefficients give
\begin{equation}
\nabla_\theta\mathcal L_{\mathrm{TRACE}}
=\frac1{|\mathcal I_+|}\sum_{j\in\mathcal I_+}
\sum_{R\in\mathcal R_j^+}\alpha_{R\mid j}
\nabla_\theta\overline H_R(\theta).
\end{equation}
A token in region $R$ has explicit entropy coefficient
$\alpha_{R\mid j_R}/(|\mathcal I_+|J|\Omega_R|)$.
Region averaging removes a direct area multiplier on the region's total coefficient; it does not remove area-dependent features, candidate selection, or shared-parameter interactions. Across overlapping crops, distinct proposals can also describe overlapping source pixels.

Equal source budgets do not imply batch invariance. Gradients still interact through common visual parameters, and Adam applied to a joint objective is not equivalent to independent per-image updates. Changing episode composition changes those interactions even when all sources retain equal coefficient mass. Neither this algebra nor a high C response establishes that an update improves unknown ground-truth labels.

\section{Global Mode Selection and Training-Source Confirmation}
\label{app:selection}
\subsection{Cohorts and Locked Selection Rule}
One global probe mode is selected using PC60 and COCO Object training-source calibration only. For each dataset, a label-free hash ranking with selection seed 20261029 partitions 128 calibration images and 256 distinct confirmation images, after excluding identifiable historical development IDs and official evaluation identities. Confirmation uses the original sorted order and two permutations with seeds 20261030 and 20261031. Model seed 0 and inherited dataset settings remain unchanged. The three orders reuse the same images and are not independent replicates.

The six arms are episode-normalized minimum (\emph{Old min}), image-normalized Uniform, Confidence, T, C, and minimum. In this comparison, $C_R$ denotes exactly the support $\rho_R$ in Eq.~\ref{eq:app-C}, not an additional score. T measures the original query after replacing context with replicated grayscale followed by $9\times9$ stride-one average blur with padding four and the library's default padded averaging; grayscale weights are $(0.2989,0.5870,0.1140)$. Candidate pixels remain unchanged. Calling that response $T_R$, the minimum alternatives use $\min(T_R,C_R)$. Only Old min uses episode normalization. T and C require one diagnostic input per candidate; either minimum arm requires both. These are fixed alternatives for selection/comparison, not components of the selected C method.

For mode $m\in\{T,C,\min\}$ and calibration dataset $d$, let $\Delta_{d,m}$ be mIoU minus image Confidence. The locked rule maximizes $\min_d\Delta_{d,m}$. Modes within 0.02 percentage points of the best worst-dataset margin are tied; ties prefer fewer actual calibration diagnostic encoder inputs, then larger unweighted mean margin, then the fixed order T, C, minimum. Both datasets and all calibration arms are prediction-sealed before calibration scoring. The selected mode is sealed before any confirmation prediction, and all confirmation arms/orders are sealed before confirmation scoring.

\subsection{Selection Outcome and Scope}
C is selected despite a negative calibration margin on PC60. Its calibration differences from image Confidence are $-0.272957$ pp on PC60 and $+0.579561$ pp on Object; its worst-dataset margin exceeds T's $-0.518667$ and minimum's $-0.308598$. The latter difference is outside the 0.02-pp tie band. Confirmation does not reselect the mode.

\begin{table*}[t]
\caption{Global C was sealed from calibration before confirmation; no reselection. Prospective training-source cohorts under partial documented history only; July candidate pools remain uncertified. Not independent-unseen or full-split evidence. All orders are dependent reorderings, not confidence intervals or significance evidence. Pooled mIoU includes warmup. Calibration128 is separate from confirmation256; all six arms and three confirmation orders are retained.}
\centering\scriptsize
\begin{tabular}{@{}lllrrrrrr@{}}
\toprule
Dataset & Stage & Order & Old min & Image uniform & Image confidence & Image T & Image C & Image min \\
\midrule
PC60 & Calibration & Sorted & 24.846289 & 25.250079 & 25.212533 & 24.693866 & 24.939576 & 24.903935 \\
Object & Calibration & Sorted & 25.799302 & 25.620627 & 25.639792 & 26.711807 & 26.219353 & 25.574809 \\
PC60 & Confirmation & Sorted & 26.569958 & 26.352055 & 26.341809 & 26.287044 & 26.518186 & 26.458232 \\
PC60 & Confirmation & 20261030 & 26.629852 & 26.496603 & 26.386787 & 26.412501 & 26.455814 & 26.456690 \\
PC60 & Confirmation & 20261031 & 26.681803 & 26.358987 & 26.390180 & 26.301490 & 26.544012 & 26.356117 \\
Object & Confirmation & Sorted & 29.955178 & 26.573758 & 27.221342 & 31.100526 & 29.324619 & 29.975437 \\
Object & Confirmation & 20261030 & 30.351900 & 26.833645 & 28.066020 & 30.881844 & 30.166647 & 30.837342 \\
Object & Confirmation & 20261031 & 29.796538 & 26.328251 & 27.041885 & 30.440226 & 29.966576 & 29.832325 \\
\bottomrule
\end{tabular}
\end{table*}
\begin{table*}[t]
\caption{Signed differences versus image confidence in percentage points. Retain negative, zero and positive differences; no post-confirmation mode choice.}
\centering\scriptsize
\begin{tabular}{@{}lllrrrrrr@{}}
\toprule
Dataset & Stage & Order & Old min & Image uniform & Image confidence & Image T & Image C & Image min \\
\midrule
PC60 & Calibration & Sorted & -0.366244 & +0.037546 & +0.000000 & -0.518667 & -0.272957 & -0.308598 \\
Object & Calibration & Sorted & +0.159510 & -0.019165 & +0.000000 & +1.072014 & +0.579561 & -0.064984 \\
PC60 & Confirmation & Sorted & +0.228150 & +0.010246 & +0.000000 & -0.054764 & +0.176377 & +0.116423 \\
PC60 & Confirmation & 20261030 & +0.243065 & +0.109815 & +0.000000 & +0.025714 & +0.069026 & +0.069902 \\
PC60 & Confirmation & 20261031 & +0.291623 & -0.031193 & +0.000000 & -0.088690 & +0.153832 & -0.034063 \\
Object & Confirmation & Sorted & +2.733836 & -0.647584 & +0.000000 & +3.879184 & +2.103277 & +2.754095 \\
Object & Confirmation & 20261030 & +2.285879 & -1.232375 & +0.000000 & +2.815823 & +2.100626 & +2.771322 \\
Object & Confirmation & 20261031 & +2.754653 & -0.713634 & +0.000000 & +3.398341 & +2.924691 & +2.790440 \\
\bottomrule
\end{tabular}
\end{table*}


Selected C exceeds image Confidence in all three confirmation orders on both datasets: $+0.069026$ to $+0.176377$ pp on PC60 and $+2.100626$ to $+2.924691$ pp on Object. This is not best-arm superiority. Old min is higher than C in all three PC60 orders, and T is higher than C in all three Object orders. The tables retain every arm and signed difference, including negative results.

The cohort evidence is prospective only relative to the documented exclusion inventory. Exact membership of historical July 4,096-image candidate pools remains uncertified because surviving loader/train-split records do not establish the necessary mapping. Neither global never-seen status nor exhaustive historical exclusion is established. The cohort source is \path{trace_image_budget_20261029_0233}, including \texttt{SELECTION\_SEALED.json}, under the results directory above. The selection seal is pinned by the new full suite. Neither these dependent reorderings nor the full seed-0 run supplies independent confidence intervals, seed dispersion, bootstrap inference, or significance.

The selected C is retained unchanged for the six full evaluation splits. Its new mIoUs are 30.693152 (PC60), 32.594583 (Object), 52.621491 (VOC21), 83.487219 (VOC20), 74.982118 (CD), and 21.836221 (Stuff). It exceeds the matched image Confidence arm on all six, with differences from $+0.042865$ pp (Stuff) to $+7.317811$ pp (CD). These are descriptive finite-panel comparisons of the fixed variant, not universal accuracy claims or a license to choose a per-dataset mode.

\section{Computation and Timing Scope}
\label{app:cost}
The selected C method independently constructs candidates, evaluates one context view per selected candidate, adapts, and reconstructs predictions. Uniform and Confidence omit diagnostic forwards but retain candidate construction and optimization. For $P$ input crops, $S$ executed update steps, and $n$ selected candidates in a nonempty episode, the standard encoder-input schedules are $(S+2)P+n$ for C and $(S+2)P$ for its zero-view controls: one initial candidate pass, $S$ adaptation passes, and one final inference pass. T has the same diagnostic-input count as C; a two-view minimum adds $2n$ instead. Actual counters, rather than this nonempty-episode formula, govern empty or skipped episodes.

Diagnostic inputs and batched encoder calls are different costs. A one-probe design halves diagnostic inputs relative to a two-probe design with identical candidates, not total model work or wall time. The full suite contains only Uniform, Confidence, and C, so it provides no full-split timing measurement for a minimum arm. No historical cohort timing ratio is transferred to these full results.

Full-split costs use synchronized model-side timers around reset, candidate construction, requested diagnostics, adaptation, final reconstruction, and CPU prediction transfer. They include the full worker's observation-only finite-value checks and accounting hooks; these checks abort on nonfinite logits, gradients, parameters, or Adam state without changing the objective, adding encoder passes, or certifying semantic safety. Hook setup and artifact serialization are outside the timer. Model loading, RGB decode, H2D/geometry preparation, disk writes, and offline scoring are excluded. This is not end-to-end throughput. Episode 0 is omitted from timing only; all its predictions, accuracy counts, and encoder work are retained. The timing denominator is therefore total source images minus the first episode, whereas work counts use all source images.

\begin{table*}[t]
\caption{Synchronized model-side ms/image for the new full-split runs. $N_t$ excludes only episode 0 per arm. Includes reset, initial readout, candidates, requested views, LayerNorm optimization, final inference and CPU prediction transfer. Excludes loading, RGB decode, H2D/geometry preparation, provenance checks, disk writes and scoring; not end-to-end throughput.}
\label{tab:image-full-cost}
\centering\small
\begin{tabular}{@{}lrrrr@{}}
\toprule
Dataset & $N_t$ & Uniform & Confidence & TRACE(C) \\
\midrule
PC60 & 5096 & 83.65 & 83.73 & 119.41 \\
Object & 4986 & 237.34 & 237.41 & 292.00 \\
VOC21 & 1433 & 278.48 & 278.24 & 330.59 \\
VOC20 & 1448 & 279.88 & 279.94 & 295.56 \\
COCO-CD & 472 & 23.38 & 23.38 & 29.86 \\
Stuff & 4995 & 129.44 & 129.35 & 186.14 \\
\bottomrule
\end{tabular}
\end{table*}
\begin{table*}[t]
\caption{Actual diagnostic encoder calls / inputs, including warmup. Uniform and Confidence compute zero probes; TRACE(C) computes one input per candidate. Counts are from these full populations, not earlier cohort ratios. Total encoder work is retained in provenance.}
\label{tab:image-full-work}
\centering\small
\begin{tabular}{@{}lrrrrr@{}}
\toprule
Dataset & $N$ & Candidates & Uniform & Confidence & TRACE(C) \\
\midrule
PC60 & 5105 & 111514 & 0 / 0 & 0 / 0 & 7317 / 111514 \\
Object & 5000 & 168669 & 0 / 0 & 0 / 0 & 10662 / 168669 \\
VOC21 & 1449 & 46892 & 0 / 0 & 0 / 0 & 2956 / 46892 \\
VOC20 & 1449 & 11583 & 0 / 0 & 0 / 0 & 1449 / 11583 \\
COCO-CD & 500 & 1998 & 0 / 0 & 0 / 0 & 125 / 1998 \\
Stuff & 5000 & 164543 & 0 / 0 & 0 / 0 & 10964 / 164543 \\
\bottomrule
\end{tabular}
\end{table*}


The calibration/confirmation cost tables below retain all six independently executed arms and every order on their own cohorts. Their timing scope and warmup convention match the definitions above, but they are separate measurements from the guarded full-split worker. They do not establish equal-budget comparisons with external methods.

\begin{table*}[t]
\caption{Global C was sealed from calibration before confirmation; no reselection. Prospective training-source cohorts under partial documented history only; July candidate pools remain uncertified. Not independent-unseen or full-split evidence. All orders are dependent reorderings, not confidence intervals or significance evidence. Synchronized model-side ms/image; $N$ excludes only episode 0. Includes reset, candidates, requested views, adaptation, final inference and CPU prediction transfer. Excludes model loading, RGB decode, H2D/geometry preparation, disk writes and scoring; not end-to-end throughput.}
\centering\scriptsize
\begin{tabular}{@{}lllrrrrrrr@{}}
\toprule
Dataset & Stage & Order & $N$ & Old min & Image uniform & Image confidence & Image T & Image C & Image min \\
\midrule
PC60 & Calibration & Sorted & 119 & 142.189 & 69.796 & 69.663 & 106.605 & 106.839 & 142.796 \\
Object & Calibration & Sorted & 114 & 327.658 & 216.959 & 216.465 & 270.054 & 271.848 & 326.164 \\
PC60 & Confirmation & Sorted & 247 & 136.390 & 66.483 & 66.448 & 101.829 & 102.030 & 136.874 \\
PC60 & Confirmation & 20261030 & 247 & 136.450 & 66.446 & 66.487 & 101.851 & 102.141 & 136.939 \\
PC60 & Confirmation & 20261031 & 247 & 136.319 & 66.588 & 66.424 & 101.655 & 101.789 & 136.771 \\
Object & Confirmation & Sorted & 242 & 314.542 & 209.470 & 209.574 & 261.083 & 264.232 & 316.154 \\
Object & Confirmation & 20261030 & 242 & 314.703 & 209.650 & 209.782 & 261.031 & 264.567 & 315.874 \\
Object & Confirmation & 20261031 & 242 & 315.985 & 210.886 & 210.439 & 261.734 & 265.267 & 317.099 \\
\bottomrule
\end{tabular}
\end{table*}
\begin{table*}[t]
\caption{Actual diagnostic encoder inputs / total encoder inputs per source image. $N$ includes all images including warmup, unlike timing. Calls and all totals are retained in provenance.}
\centering\scriptsize
\begin{tabular}{@{}lllrrrrrrr@{}}
\toprule
Dataset & Stage & Order & $N$ & Old min & Image uniform & Image confidence & Image T & Image C & Image min \\
\midrule
PC60 & Calibration & Sorted & 128 & 43.688 / 56.688 & 0.000 / 13.000 & 0.000 / 13.000 & 21.844 / 34.844 & 21.844 / 34.844 & 43.688 / 56.688 \\
Object & Calibration & Sorted & 128 & 67.609 / 98.609 & 0.000 / 31.000 & 0.000 / 31.000 & 33.805 / 64.805 & 33.805 / 64.805 & 67.609 / 98.609 \\
PC60 & Confirmation & Sorted & 256 & 43.703 / 56.703 & 0.000 / 13.000 & 0.000 / 13.000 & 21.852 / 34.852 & 21.852 / 34.852 & 43.703 / 56.703 \\
PC60 & Confirmation & 20261030 & 256 & 43.703 / 56.703 & 0.000 / 13.000 & 0.000 / 13.000 & 21.852 / 34.852 & 21.852 / 34.852 & 43.703 / 56.703 \\
PC60 & Confirmation & 20261031 & 256 & 43.703 / 56.703 & 0.000 / 13.000 & 0.000 / 13.000 & 21.852 / 34.852 & 21.852 / 34.852 & 43.703 / 56.703 \\
Object & Confirmation & Sorted & 256 & 67.680 / 98.680 & 0.000 / 31.000 & 0.000 / 31.000 & 33.840 / 64.840 & 33.840 / 64.840 & 67.680 / 98.680 \\
Object & Confirmation & 20261030 & 256 & 67.680 / 98.680 & 0.000 / 31.000 & 0.000 / 31.000 & 33.840 / 64.840 & 33.840 / 64.840 & 67.680 / 98.680 \\
Object & Confirmation & 20261031 & 256 & 67.680 / 98.680 & 0.000 / 31.000 & 0.000 / 31.000 & 33.840 / 64.840 & 33.840 / 64.840 & 67.680 / 98.680 \\
\bottomrule
\end{tabular}
\end{table*}


\section{Qualitative Evidence and Selection Limits}
\label{app:figures}
A candidate illustration is an input-space view paired with a stored token-space measurement. The full C artifacts retain candidate query/category identity, source/crop IDs, confidence, measured C, detached objective coefficient, and lossless original token masks. Candidate-preserving context views can be rendered from these masks and source RGB inputs without a new diagnostic inference. Display names denote predicted queries or their recorded category mapping, not ground-truth identification. The original token footprint, not a newly segmented region in the transformed image, defines the displayed support.

Unmeasured T is unavailable in the C full suite, not zero. Whole-image pre-update category masks are also unavailable: no additional pre-update inference or independently validated reconstruction of initial logits was saved. An original candidate's mapped category is not a complete pre-update segmentation. The full artifacts therefore support single-probe candidate displays and comparisons among saved final Uniform/Confidence/C segmentations, not two-probe complementarity plots or whole-image before/after transition rates.

The teaser selects two support extremes without ground truth. Within the first 400 images of each frozen dataset order, eligible candidates are non-background predictions with $\kappa_R\geq0.9$ and 16--192 native tokens. The highest C response in PC60 and lowest in Object are selected, breaking ties by source-image index then receipt candidate index. They are PC60 image 151, candidate 158 (sofa; confidence 0.984197, support 0.997586), and Object image 241, candidate 108 (cake; confidence 0.983545, support 0.088384). Image indices are zero-based in the frozen full panels; candidate indices are episode-receipt row indices. Source files are \texttt{2009\_001581.jpg} and \texttt{000000025096.jpg}, respectively. The rendered context views use the exact float32 transform and original $16\times16$ masks; displayed supports are saved measurements, not new model evaluations.

Segmentation cases are offline named-class IoU extremes among the first 400 images per dataset. For each image, choose its largest non-background ground-truth category with at least 128 valid pixels, breaking area ties by the lowest category ID. Compute that single category's IoU for saved Confidence and C predictions. Select the largest C-minus-Confidence difference for CD, Object, and PC60, and the smallest for VOC21; ties use the lowest image index. The selected cases are CD 394 (cat; 0.00 to 56.16), Object 40 (dining table; 1.19 to 86.80), PC60 199 (bus; 44.72 to 78.06), and VOC21 85 (table; 84.34 to 62.87). These values are per-image named-class IoUs, not image-averaged or dataset mIoU. The Confidence column is the matched allocation control, not MLMP.

Selected extremes illustrate both directions of change, not their frequency. Ground truth enters segmentation-display selection only after the global prediction seal and never enters adaptation or mode selection. The displays cannot establish representative prevalence, calibrated reliability, or a causal effect of one region's coefficient under a joint update. Exact identities, selection rules, stored measurements, rendering checks, and source seals are in \path{figure/image_full_provenance.json}. Earlier two-view figures and their scores are not evidence for the new C method.

\section{Reported and Local Reference Protocols}
\label{app:native}
Reported operating points (R) and local measurements (L) have different evidential roles. For VOC20, VOC21, PC60, COCO Object, and COCO-Stuff, MLMP's main Table 5 reports clean NACLIP ViT-L/14 values 83.76, 50.78, 27.99, 28.84, and 21.25~\cite{noori2025mlmp}. The new local C values exceed four of those five references, but VOC20 is 83.487219, below the reported 83.76 by 0.272781 pp. We do not claim that all reported references are surpassed. A shared backbone family does not certify identical checkpoints, query vocabularies, loaders, readouts, optimization budgets, or tuning histories.

COCO-CD has a separate completed local native reference. The adapting native baselines retain their backend behavior and prompts with learning rate 0.001 and 10 updates, whereas C uses Table~\ref{tab:configurations}. The best retained native baseline is TENT~\cite{wang2021tent} at 72.200299, compared with new C at 74.982118. Although these runs share the ordered input panel and episode composition, their native protocols differ; the 2.781819-pp difference is not an equal-budget attribution of the probe or the objective. The historical native No Adapt row is not a certified reset readout of the new method and does not establish its before/after adaptation gain.

Published statistics remain attributed to their sources and are not pooled with local deterministic results. Existing published MLMP values are used without additional MLMP reruns. No local WATT or CLIPArTT evaluation is claimed. The within-implementation Uniform/Confidence/C controls, training-source selection comparisons, cross-paper operating points, and local native CD reference answer different questions and are not interchangeable.


\end{document}
